% ===== Topic 1, part d (EN): problems 41-50 (challenge problems extending problems 16-30) =====
% helper shapes for problem 41 (own prefix \rvab..., independent of part b)
\providecommand{\rvabShp}[3]{%
  \if#1T\draw[tmBlue!80!black,line width=0.8pt,fill=tmBlue!18] (#2,#3)--(#2+0.34,#3)--(#2+0.17,#3+0.3)--cycle;\fi
  \if#1C\draw[tmOrange!85!black,line width=0.8pt,fill=tmOrange!22] (#2+0.15,#3+0.15) circle (0.15);\fi
  \if#1S\draw[tmGreen!80!black,line width=0.8pt,fill=tmGreen!18] (#2,#3) rectangle (#2+0.3,#3+0.3);\fi
  \if#1D\draw[tmPurple!85!black,line width=0.8pt,fill=tmPurple!18] (#2+0.13,#3)--(#2+0.26,#3+0.16)--(#2+0.13,#3+0.32)--(#2,#3+0.16)--cycle;\fi
  \if#1Q\node[font=\Large\bfseries,text=Ink,inner sep=0pt,anchor=south] at (#2+0.15,#3) {?};\fi}
\providecommand{\rvabTri}{\tikz[baseline=0pt]{\rvabShp{T}{0}{0}}}
\providecommand{\rvabCir}{\tikz[baseline=0pt]{\rvabShp{C}{0}{0}}}
\providecommand{\rvabSq}{\tikz[baseline=0pt]{\rvabShp{S}{0}{0}}}
\providecommand{\rvabDia}{\tikz[baseline=0pt]{\rvabShp{D}{0}{0}}}
\providecommand{\rvabRow}[3]{%
  \foreach \rvabS [count=\rvabI] in {#1}{\xdef\rvabN{\rvabI}}%
  \foreach \rvabS [count=\rvabI] in {#1}{\rvabShp{\rvabS}{#2-0.19*\rvabN+0.38*\rvabI-0.36}{#3+0.03}}}
\providecommand{\rvabScale}[4]{\begin{scope}[shift={(#1,#2)}]
  \draw[Ink,line width=1pt] (-1.9,0.55)--(-0.5,0.55) (-1.2,0.55)--(-1.2,0.25)--(1.2,0.25)--(1.2,0.55) (0.5,0.55)--(1.9,0.55);
  \draw[Ink,line width=0.9pt,fill=Ink!12] (-0.2,-0.12)--(0,0.25)--(0.2,-0.12)--cycle;
  \rvabRow{#3}{-1.2}{0.55}\rvabRow{#4}{1.2}{0.55}
\end{scope}}

% ------------------------------------------------------------------ 41
\begin{probbox}{41}{\stars{1}}
The three scales below are balanced; identical shapes have equal masses.
\begin{figC}
\rvabScale{0}{0}{T,T}{C,C,C}
\rvabScale{4.4}{0}{C,S}{T}
\rvabScale{8.8}{0}{D}{T,S}
\node[font=\footnotesize\bfseries,text=Slate] at (0,-0.45) {Scale 1};
\node[font=\footnotesize\bfseries,text=Slate] at (4.4,-0.45) {Scale 2};
\node[font=\footnotesize\bfseries,text=Slate] at (8.8,-0.45) {Scale 3};
\end{figC}
\begin{enumerate}[label=\alph*),leftmargin=1.6em,itemsep=1pt]
\item How many \rvabSq\ shapes balance one \rvabCir?
\item How many \rvabSq\ shapes balance one \rvabDia?
\item One \rvabDia\ and one \rvabCir\ are placed on the left pan. If only \rvabTri\ shapes may be used on the right pan, how many \rvabTri\ are needed?
\item The four objects \rvabTri, \rvabCir, \rvabSq, \rvabDia\ (one of each kind) have a total mass of $50$~g. Find the mass of each object.
\end{enumerate}
\end{probbox}
\probsol{
\textbf{Analysis.} This extends Problem 17. Let $T$, $C$, $S$, $D$ (grams) be the masses of one \rvabTri, one \rvabCir, one \rvabSq\ and one \rvabDia. Each balanced scale gives an equation (mass on the left pan $=$ mass on the right pan):
\[ \text{(1) } 2T=3C,\qquad \text{(2) } C+S=T,\qquad \text{(3) } D=T+S. \]
Idea: express every mass in terms of the same unknown, $S$.
\begin{recallbox}Substitution: if we know $T=C+S$, we may replace $T$ by $C+S$ in any equation. Adding or subtracting the same expression on both sides of an equation gives a true equation.\end{recallbox}

\step{1}{Part a -- Substitute (2) into (1)}
Replace $T$ by $C+S$ in (1): $2(C+S)=3C$, that is $2C+2S=3C$.\par
Subtract $2C$ from both sides: $2S=3C-2C=C$. So $C=2S$.\par
One \rvabCir\ weighs the same as two \rvabSq: we need \textbf{$2$ shapes \rvabSq}.
\begin{tipbox}Why substitute (2) into (1)? Because (1) contains only $T$ and $C$; after replacing $T=C+S$, the equation contains only $C$ and $S$, so we can compare $C$ with $S$ at once.\end{tipbox}

\step{2}{Express $T$ in terms of $S$}
From (2): $T=C+S=2S+S=3S$.\par
(Check with (1): $2T=6S$ and $3C=3\cdot 2S=6S$ \checkmark.)

\step{3}{Part b}
From (3): $D=T+S=3S+S=4S$. We need \textbf{$4$ shapes \rvabSq}.

\step{4}{Part c}
Left pan: $D+C=4S+2S=6S$. Each \rvabTri\ weighs $T=3S$, so the number of \rvabTri\ needed is $6S\div 3S=$ \textbf{$2$ shapes \rvabTri}. (Indeed, $2T=6S=D+C$.)

\step{5}{Part d -- An equation in one unknown}
$T+C+S+D=50$; in terms of $S$: $3S+2S+S+4S=50$, that is $10S=50$, so $S=5$. Hence:
\begin{center}
\begin{tabular}{|c|c|c|c|c|}
\hline
\rowcolor{tmBlue!10}Object & \rvabSq & \rvabCir & \rvabTri & \rvabDia\\ \hline
In terms of $S$ & $S$ & $C=2S$ & $T=3S$ & $D=4S$\\ \hline
Mass & $5$~g & $10$~g & $15$~g & $20$~g\\ \hline
\end{tabular}
\end{center}
\verify{Scale 1: $2\cdot 15=30=3\cdot 10$ \checkmark. Scale 2: $10+5=15$ \checkmark. Scale 3: $20=15+5$ \checkmark. Total: $15+10+5+20=50$ \checkmark.}
\begin{parentbox}Ask your child to read each scale as an equation with letters ($2T=3C$, \dots) before calculating. ``Removing the same object from both pans'' is exactly subtracting the same number from both sides -- the foundation of solving equations.\end{parentbox}
\ans{a) $2$ shapes \rvabSq;\ b) $4$ shapes \rvabSq;\ c) $2$ shapes \rvabTri;\ d) \rvabSq\ $5$~g, \rvabCir\ $10$~g, \rvabTri\ $15$~g, \rvabDia\ $20$~g.}
}

% ------------------------------------------------------------------ 42
\begin{probbox}{42}{\stars{1}}
In this problem we only consider triangles in which every angle measures (in degrees) a \emph{whole number}. Two isosceles triangles are said to be of the \emph{same type} if they have the same three angle measures.
\begin{enumerate}[label=\alph*),leftmargin=1.6em,itemsep=1pt]
\item Prove that there is no triangle whose three angles all have odd measures.
\item How many types of isosceles triangle have all three angles with even measures?
\item How many types of isosceles triangle have exactly two angles with odd measures? Is there an isosceles triangle with exactly one angle of odd measure?
\item How many types of isosceles triangle are there altogether?
\end{enumerate}
\end{probbox}
\probsol{
\textbf{Analysis.} This extends Problem 20: instead of looking only at triangles with a $70^\circ$ angle, we look at \emph{all} isosceles triangles with whole-number angles and sort them by whether their angles are even or odd.
\begin{recallbox}The three angles of a triangle add up to $180^\circ$. An isosceles triangle has two equal base angles. Even and odd: odd $+$ odd $=$ even; even $+$ odd $=$ odd; the sum of three odd numbers is odd; $2\times$(any number) is even.\end{recallbox}

\step{1}{Part a}
Suppose all three angles are odd. The sum of the first two odd numbers is even; adding one more odd number gives an \emph{odd} number. But the sum of the three angles is $180$ -- an even number. Contradiction. So no triangle has three odd angles.

\step{2}{Describe every isosceles triangle by one number $x$}
Let each base angle measure $x$ (degrees). Then the apex angle is $180-2x$. For a triangle, the apex angle must be positive, $180-2x\ge 1$, that is, $x\le 89$. So $x$ can be $1,2,3,\dots,89$, and each value of $x$ gives exactly one type of isosceles triangle.
\begin{proofbox}Two different values of $x$ give two different sets of angles (the base angles differ), and every isosceles triangle has the form $(x,x,180-2x)$. The equilateral triangle $(60,60,60)$ corresponds to $x=60$ -- it is counted only once. Note: we can always take the two equal angles as the base angles; for the equilateral triangle any angle will do.\end{proofbox}

\step{3}{The apex angle is always even}
$180-2x=2\cdot(90-x)$ is an even number for every $x$. So the parity pattern of the triangle depends only on $x$:
\begin{itemize}[leftmargin=1.3em,itemsep=1pt]
\item $x$ even: two even base angles, even apex angle $\Rightarrow$ all three angles even.
\item $x$ odd: two odd base angles, even apex angle $\Rightarrow$ exactly two angles odd.
\end{itemize}

\step{4}{Part b}
Count the even numbers from $1$ to $89$: $2,4,\dots,88$, that is $88\div 2=44$ numbers. So there are \textbf{$44$ types}.

\step{5}{Part c}
Count the odd numbers from $1$ to $89$: $1,3,\dots,89$, that is $(89+1)\div 2=45$ numbers. So \textbf{$45$ types} have exactly two odd angles.\par
An isosceles triangle with exactly one odd angle does \textbf{not exist}: by Step 3, the number of odd angles can only be $0$ or $2$.

\step{6}{Part d}
Total number of types: $44+45=89$ (exactly the number of values of $x$).
\begin{center}\small
\begin{tabular}{|c|c|c|}
\hline
\rowcolor{tmBlue!10}$x$ (base angle) & The three angles & Odd angles\\ \hline
$1$ & $1^\circ,\,1^\circ,\,178^\circ$ & $2$\\ \hline
$2$ & $2^\circ,\,2^\circ,\,176^\circ$ & $0$\\ \hline
$55$ & $55^\circ,\,55^\circ,\,70^\circ$ & $2$\\ \hline
$70$ & $70^\circ,\,70^\circ,\,40^\circ$ & $0$\\ \hline
$89$ & $89^\circ,\,89^\circ,\,2^\circ$ & $2$\\ \hline
\end{tabular}
\end{center}
\begin{warnbox}Do not forget the ``very flat'' triangles such as $(1^\circ,1^\circ,178^\circ)$ -- they are still genuine triangles. And do not allow $x=90$: then the apex angle would be $0^\circ$, which is no longer a triangle.\end{warnbox}
\verify{Listing $x=1,\dots,89$ with a computer gives exactly $44$ all-even triples and $45$ triples with two odd angles. The rows $x=55$ and $x=70$ are exactly the two triangles of Problem 20.}
\ans{a) It does not exist;\ b) $44$;\ c) $45$ types, and no isosceles triangle has exactly one odd angle;\ d) $89$.}
}

% ------------------------------------------------------------------ 43
\begin{probbox}{43}{\stars{2}}
A dot starts at the origin $O(0,0)$ and moves by the rule of Problem 22: the first move is $1$ unit horizontally or vertically; each later move is $1$ unit longer than the previous one and perpendicular to the move just before it (either direction allowed).
\begin{enumerate}[label=\alph*),leftmargin=1.6em,itemsep=1pt]
\item Explain why after exactly $10$ moves the dot cannot be back at $O$.
\item Find the smallest $n$ ($n\ge 1$) such that after exactly $n$ moves the dot can be back at $O$. Show one way to do it.
\item Can the dot be back at $O$ after exactly $8$ moves?
\end{enumerate}
\end{probbox}
\probsol{
\textbf{Analysis.} As in Problem 22, the moves alternate horizontal -- vertical. One direction gets the moves of \emph{odd} length $1,3,5,\dots$; the other gets the moves of \emph{even} length $2,4,6,\dots$ The dot returns to $O$ exactly when the signed sum of \emph{each} group is $0$.
\begin{recallbox}Being back at $O$ means the total horizontal displacement is $0$ \emph{and} the total vertical displacement is $0$. Changing the sign of a term $a$ in a sum changes the sum by $2a$ -- an even number -- so it \textbf{does not change whether the sum is even or odd}.\end{recallbox}

\step{1}{Part a}
After $10$ moves, the odd group consists of the lengths $1,3,5,7,9$ (whether the first move is horizontal or vertical). The signed sum $\pm1\pm3\pm5\pm7\pm9$ has the same parity as $1+3+5+7+9=25$, so it is always \emph{odd} and can never be $0$. So in that direction the dot is not back where it started, hence the dot is not at $O$.

\step{2}{Part b -- Rule out $n\le 6$}
Look at the odd group (in one direction) for each $n$:
\begin{center}
\begin{tabular}{|c|c|c|}
\hline
\rowcolor{tmBlue!10}$n$ & Odd lengths & Why the signed sum cannot be $0$\\ \hline
$1,\ 2$ & $\{1\}$ & $\pm 1\neq 0$\\ \hline
$3,\ 4$ & $\{1,3\}$ & $\pm1\pm3$ can only be $4, 2, -2, -4$\\ \hline
$5,\ 6$ & $\{1,3,5\}$ & the signed sum has the parity of $1+3+5=9$ (odd)\\ \hline
\end{tabular}
\end{center}
So for $n\le 6$ the dot cannot return to $O$, and $n\ge 7$.

\step{3}{Part b -- $n=7$ works}
Odd group: $\{1,3,5,7\}$, choose $+1-3-5+7=0$. Even group: $\{2,4,6\}$, choose $+2+4-6=0$.\par
Let the first move be horizontal. The moves: right $1$, up $2$, left $3$, up $4$, left $5$, down $6$, right $7$. The positions are
\[ (0,0)\to(1,0)\to(1,2)\to(-2,2)\to(-2,6)\to(-7,6)\to(-7,0)\to(0,0). \]
\begin{figC}\begin{scope}[scale=0.5]
\draw[Line,line width=0.3pt] (-8,-1) grid (2,7);
\draw[->,Navy,line width=0.5pt] (-8,0)--(2.4,0); \draw[->,Navy,line width=0.5pt] (0,-1)--(0,7.4);
\draw[tmBlue,line width=1.3pt,->] (0,0)--(1,0)--(1,2)--(-2,2)--(-2,6)--(-7,6)--(-7,0)--(-0.32,0);
\fill[AnsRed] (0,0) circle (3pt);
\node[font=\scriptsize,anchor=north east,inner sep=1.5pt] at (0,0) {$O$};
\foreach \p/\l/\a in {(0.5,0)/1/below,(1,1)/2/right,(-0.5,2)/3/below,(-2,4)/4/right,(-4.5,6)/5/above,(-7,3)/6/left,(-3.5,0)/7/below}{\node[\a,font=\scriptsize\bfseries,text=tmBlue!80!black] at \p {\l};}
\end{scope}\end{figC}
So the smallest $n$ is $7$.

\step{4}{Part c -- $n=8$}
Odd group: $\{1,3,5,7\}$, again $+1-3-5+7=0$. Even group: $\{2,4,6,8\}$, choose $+2-4-6+8=0$. So \textbf{yes}: right $1$, up $2$, left $3$, down $4$, left $5$, down $6$, right $7$, up $8$; the positions are $(1,0)$, $(1,2)$, $(-2,2)$, $(-2,-2)$, $(-7,-2)$, $(-7,-8)$, $(0,-8)$, $(0,0)$.
\begin{proofbox}A signed sum $\pm a_1\pm a_2\pm\dots$ equals $0$ exactly when the numbers $a_1,a_2,\dots$ can be split into two groups with equal sums (the group with sign $+$ and the group with sign $-$). For example $\{1,3,5,7\}$: $1+7=3+5=8$. That is why the total of all the numbers must be \emph{even}.\end{proofbox}
\begin{warnbox}Do not forget that \emph{both} groups must sum to $0$. For $n=5$: the even group $\{2,4\}$ cannot give $0$, and neither can the odd group $\{1,3,5\}$ -- one failing group is already enough to rule $n$ out.\end{warnbox}
\verify{Checking every choice of directions with a computer for $n=1,2,\dots,15$: the dot can return to $O$ for $n=7$, $n=8$ and $n=15$; it cannot for any $n\le 6$.}
\ans{a) The signed sum of $1,3,5,7,9$ is always odd, so it is not $0$;\ b) $n=7$ (for example: right 1, up 2, left 3, up 4, left 5, down 6, right 7);\ c) Yes.}
}

% ------------------------------------------------------------------ 44
\begin{probbox}{44}{\stars{2}}
A $4\times 6\times n$ rectangular box ($n$ a positive integer) is built from $1\times1\times1$ unit cubes. Let $A$ be the surface area of the box and $B$ the combined surface area of all the unit cubes.
\begin{enumerate}[label=\alph*),leftmargin=1.6em,itemsep=1pt]
\item Write $A$, $B$ and the ratio $\dfrac{B}{A}$ as expressions in $n$.
\item Prove that $\dfrac{B}{A}<7.2$ for every $n$.
\item Find all values of $n$ for which $\dfrac{B}{A}$ is an integer.
\end{enumerate}
\end{probbox}
\probsol{
\textbf{Analysis.} This extends Problem 24, but this time the trick of ``splitting off the whole part'' is not as neat. We use a different method: bound the ratio from above, then test each possible integer value.
\begin{recallbox}An $a\times b\times c$ box: surface area $2(ab+bc+ca)$, made of $a\cdot b\cdot c$ unit cubes; each unit cube has surface area $6$.\end{recallbox}

\step{1}{Part a}
$A=2(4\cdot 6+4\cdot n+6\cdot n)=2(24+10n)=48+20n$.\par
Number of unit cubes: $4\cdot 6\cdot n=24n$; therefore $B=6\cdot 24n=144n$.\par
\[ \frac{B}{A}=\frac{144n}{48+20n}=\frac{4\cdot 36n}{4\cdot(12+5n)}=\frac{36n}{5n+12}. \]

\step{2}{Part b}
Since $5n+12>5n$, we are dividing $36n$ by a number \emph{larger} than $5n$, so the result is \emph{smaller}:
\[ \frac{36n}{5n+12}<\frac{36n}{5n}=\frac{36}{5}=7.2. \]
\begin{recallbox}For positive numbers: with the same numerator, the fraction with the larger denominator is smaller.\end{recallbox}

\step{3}{Part c -- Call the integer value $k$}
If $\frac{B}{A}=k$ is an integer, then $k$ is positive and $k<7.2$, so $k\in\{1,2,3,4,5,6,7\}$. We have
\[ \frac{36n}{5n+12}=k \iff 36n=k(5n+12)=5kn+12k \iff 36n-5kn=12k \iff n(36-5k)=12k. \]
So $n=\dfrac{12k}{36-5k}$, and we only need to test $7$ values of $k$.

\step{4}{Test each value of $k$}
\begin{center}
\begin{tabular}{|c|c|c|c|c|c|c|c|}
\hline
\rowcolor{tmBlue!10}$k$ & $1$ & $2$ & $3$ & $4$ & $5$ & $6$ & $7$\\ \hline
$12k$ & $12$ & $24$ & $36$ & $48$ & $60$ & $72$ & $84$\\ \hline
$36-5k$ & $31$ & $26$ & $21$ & $16$ & $11$ & $6$ & $1$\\ \hline
$n=\frac{12k}{36-5k}$ & $\frac{12}{31}$ $\times$ & $\frac{24}{26}$ $\times$ & $\frac{36}{21}$ $\times$ & $\mathbf{3}$ & $\frac{60}{11}$ $\times$ & $\mathbf{12}$ & $\mathbf{84}$\\ \hline
\end{tabular}
\end{center}
(The mark $\times$ means: not an integer.)
So $n\in\{3,\,12,\,84\}$.
\begin{proofbox}Why have we missed nothing? Every $n$ making $\frac{B}{A}$ an integer gives an integer $k$ from $1$ to $7$ (part b), and for each $k$ there is exactly one $n$ with $n(36-5k)=12k$. We have tested all $7$ possibilities.\end{proofbox}
\verify{
$n=3$: $A=48+60=108$, $B=432$, $\frac{B}{A}=4$ \checkmark.\quad
$n=12$: $A=48+240=288$, $B=1728$, $\frac{B}{A}=6$ \checkmark.\quad
$n=84$: $A=48+1680=1728$, $B=12\,096$, $\frac{B}{A}=7$ \checkmark.}
\begin{tipbox}When a fraction containing $n$ cannot be ``split'' neatly, find a number it always stays below. If only a few integer values remain possible, test them all!\end{tipbox}
\ans{a) $A=48+20n$, $B=144n$, $\dfrac{B}{A}=\dfrac{36n}{5n+12}$;\ b) see the solution;\ c) $n=3$, $n=12$, $n=84$.}
}

% ------------------------------------------------------------------ 45
\begin{probbox}{45}{\stars{2}}
Two bus routes cross at the stop $X$. Bus A runs back and forth along the route $P\text{--}X\text{--}S$; each arm ($PX$ or $XS$) takes exactly $4$ minutes. Bus B runs along the route $Q\text{--}X\text{--}T$; each arm takes exactly $6$ minutes. At 1:00 p.m., bus A leaves $P$ and bus B leaves $Q$; the buses turn around instantly at the ends of their routes and run continuously.
\begin{enumerate}[label=\alph*),leftmargin=1.6em,itemsep=1pt]
\item Prove that buses A and B are never at $X$ at the same time.
\item A bus C is added on the route $R\text{--}X\text{--}U$; each arm takes $10$ minutes, and it leaves $R$ at 1:00 p.m. Between 1:00 p.m. and 7:00 p.m., at which times are two buses at $X$ together?
\item Suppose each arm of bus C's route takes $m$ minutes ($m$ any positive integer). Prove that bus C cannot both meet bus A at $X$ (at some moment) and meet bus B at $X$ (at some moment).
\end{enumerate}
\end{probbox}
\probsol{
\textbf{Analysis.} This extends Problem 25. Measure time $t$ in minutes after 1:00 p.m. As in Problem 25, a bus whose arm takes $a$ minutes first reaches $X$ at minute $a$, and then returns to $X$ every $2a$ minutes (out to the end and back):
\[ t=a,\ 3a,\ 5a,\ 7a,\dots \quad\text{that is } t=a\times(\text{odd number}). \]
So bus A is at $X$ when $t=4\times(\text{odd})$; bus B when $t=6\times(\text{odd})$; bus C when $t=m\times(\text{odd})$.
\begin{recallbox}The product of two odd numbers is odd. A number divisible by $4$ gives an even number when divided by $2$.\end{recallbox}

\step{1}{Part a}
Times when bus A is at $X$: $t=4\times(\text{odd})$ -- always \textbf{divisible by $4$}.\par
Times when bus B is at $X$: $t=6\times(\text{odd})=2\times(3\times\text{odd})=2\times(\text{odd})$ -- dividing by $2$ gives an odd number, so it is \textbf{not divisible by $4$}.\par
A number cannot be both divisible and not divisible by $4$, so A and B are never at $X$ together.

\step{2}{Part b -- bus C with $m=10$}
Bus C is at $X$ when $t=10\times(\text{odd})=2\times(5\times\text{odd})$ -- not divisible by $4$, so C \textbf{never} meets A at $X$ (same reasoning as part a).\par
C meets B when $t=6\times(\text{odd})=10\times(\text{odd})$. Then $\frac{t}{2}=3\times(\text{odd})=5\times(\text{odd})$ is an odd number divisible by $3$ and $5$, i.e.\ an odd multiple of $15$. So $t=30\times(\text{odd})$. Conversely, $t=30\times(\text{odd})=6\times(5\times\text{odd})=10\times(3\times\text{odd})$ is indeed a time when both buses are at $X$.\par
From 1:00 p.m. to 7:00 p.m. means $0\le t\le 360$: $t=30,\ 90,\ 150,\ 210,\ 270,\ 330$.
\begin{center}
\begin{tabular}{|c|c|c|c|c|c|c|}
\hline
\rowcolor{tmBlue!10}$t$ (min) & $30$ & $90$ & $150$ & $210$ & $270$ & $330$\\ \hline
Time & 1:30 p.m. & 2:30 p.m. & 3:30 p.m. & 4:30 p.m. & 5:30 p.m. & 6:30 p.m.\\ \hline
\end{tabular}
\end{center}
At all $6$ of these times the two buses at $X$ are B and C.

\step{3}{Part c -- ``the number of halvings''}
For each positive integer $t$, keep dividing $t$ by $2$ until the result is odd. Let $h(t)$ be the number of divisions. For example: $12\to 6\to 3$, so $h(12)=2$; $6\to3$, so $h(6)=1$; $5$ is already odd, so $h(5)=0$.
\begin{proofbox}\textbf{Multiplying by an odd number does not change $h$.} Write $m=2\cdot2\cdots2\times r$ (with $h(m)$ factors $2$ and $r$ odd). Then $m\times(\text{odd})=2\cdot2\cdots2\times(r\times\text{odd})$, where $r\times\text{odd}$ is still odd. So $m\times(\text{odd})$ can also be halved exactly $h(m)$ times.\end{proofbox}
Applying this:
\begin{itemize}[leftmargin=1.3em,itemsep=1pt]
\item Every time bus A is at $X$: $h\big(4\times\text{odd}\big)=h(4)=2$.
\item Every time bus B is at $X$: $h\big(6\times\text{odd}\big)=h(6)=1$.
\item Every time bus C is at $X$: $h\big(m\times\text{odd}\big)=h(m)$.
\end{itemize}
If C meets A then $h(m)=2$; if C meets B then $h(m)=1$. A number $m$ cannot have both $h(m)=2$ and $h(m)=1$. So C cannot meet both A and B.
\begin{warnbox}Part c does not say that C \emph{never} meets anyone: with $m=4$, C meets A (at $t=4$); with $m=2$, C meets B (at $t=6$). What we proved is that C cannot meet \emph{both}.\end{warnbox}
\verify{Checking with a computer for $m=1,2,\dots,12$: no value of $m$ lets C meet both A and B; $m=4$ and $m=12$ meet A; $m=2,6,10$ meet B.}
\ans{a) Times when A is at $X$ are divisible by $4$, times when B is at $X$ are not;\ b) 1:30, 2:30, 3:30, 4:30, 5:30 and 6:30 p.m. (always buses B and C);\ c) see the solution.}
}

% ------------------------------------------------------------------ 46
\begin{probbox}{46}{\stars{1}}
Sam counts forward $7, 10, 13,\dots$ (adding $3$ each time) and stops before going past $2026$. Tom counts backward $2026, 2021, 2016,\dots$ (subtracting $5$ each time) and stops before going below $7$.
\begin{enumerate}[label=\alph*),leftmargin=1.6em,itemsep=1pt]
\item How many numbers appear in both lists?
\item Find the sum of all the numbers in part a.
\item Zoe writes the list of positive multiples of $4$: $4, 8, 12,\dots$ How many numbers appear in all three lists (Sam's, Tom's and Zoe's)? What is the largest of them?
\end{enumerate}
\end{probbox}
\probsol{
\textbf{Analysis.} This extends Problem 26. As in Problem 26: a number $N$ is in Sam's list $\iff$ it leaves remainder $1$ on division by $3$ and $7\le N\le 2026$ (because $7=3\cdot 2+1$, and adding $3$ does not change the remainder on division by $3$); a number $N$ is in Tom's list $\iff$ it leaves remainder $1$ on division by $5$ and $7\le N\le 2026$ (because $2026=5\cdot 405+1$). Now we find \emph{all} the common numbers.

\step{1}{Part a -- The test for a common number}
$N$ is common $\iff$ $N$ leaves remainder $1$ on division by $3$ and by $5$ $\iff$ $N-1$ is divisible by $3$ and by $5$ $\iff$ $N-1$ is divisible by $15$. So $N=15j+1$.
\begin{recallbox}A number divisible by $3$ and by $5$ (two coprime numbers) is divisible by $3\cdot 5=15$.\end{recallbox}

\step{2}{Part a -- Count}
We need $7\le 15j+1\le 2026$. Smallest: $j=1$ gives $16$ ($j=0$ gives $1<7$, rejected). Largest: $15j+1\le 2026 \Rightarrow 15j\le 2025 \Rightarrow j\le 135$; $j=135$ gives $2026$.\par
So the common numbers are $16, 31, 46,\dots, 2026$, for $j=1,2,\dots,135$: \textbf{$135$ numbers}.

\step{3}{Part b -- The sum (Gauss's pairing method)}
\begin{recallbox}Sum of an equally spaced list $=$ (first $+$ last) $\times$ (number of terms) $\div 2$. Reason: write the list forwards and backwards and add in pairs; every pair equals first $+$ last.\end{recallbox}
\[ 16+31+\dots+2026=\frac{(16+2026)\cdot 135}{2}=\frac{2042\cdot 135}{2}=1021\cdot 135=137\,835. \]

\step{4}{Part c -- Add the condition ``divisible by $4$''}
Look at the common numbers $16, 31, 46, 61, 76, 91, 106, 121, 136,\dots$ Which are divisible by $4$?
\[ 16\ \checkmark,\ 31,\ 46,\ 61,\ 76\ \checkmark,\ 91,\ 106,\ 121,\ 136\ \checkmark,\dots \]
In every $4$ consecutive terms exactly $1$ is divisible by $4$, and these numbers are $4\cdot 15=60$ apart: $16, 76, 136, \dots$
\begin{proofbox}Each step along the list adds $15$; adding $15$ increases the remainder on division by $4$ by $3$ (because $15=12+3$). The remainders of $16,31,46,61$ on division by $4$ are $0,3,2,1$, and then they repeat. So exactly one term in every $4$ is divisible by $4$, and two consecutive such terms are $4\times15=60$ apart.\end{proofbox}
The numbers common to all three lists: $16+60i$ with $16+60i\le 2026$, i.e.\ $60i\le 2010$, $i\le 33.5$, so $i=0,1,\dots,33$: \textbf{$34$ numbers}. The largest is $16+60\cdot 33=16+1980=1996$.
\verify{$1996=4\cdot 499$ \checkmark; $1996-7=1989=3\cdot 663$ \checkmark; $2026-1996=30=5\cdot 6$ \checkmark. The next one, $1996+60=2056>2026$, is rejected.}
\ans{a) $135$ numbers;\ b) $137\,835$;\ c) $34$ numbers, the largest is $1996$.}
}

% ------------------------------------------------------------------ 47
\begin{probbox}{47}{\stars{2}}
Emma and Noah both start at point $A$ on a circular track; each runs at a constant speed.
\begin{enumerate}[label=\alph*),leftmargin=1.6em,itemsep=1pt]
\item They run in opposite directions, with the ratio of speeds Emma : Noah equal to $5:3$. At how many different points of the track do they meet?
\item If they run in the \emph{same direction} (still with ratio $5:3$), at how many different points does Emma overtake Noah?
\item Find all speed ratios $p:q$ (with $p>q$, where $p$ and $q$ are coprime positive integers) such that, running in opposite directions, they meet at exactly $10$ different points.
\end{enumerate}
\end{probbox}
\probsol{
\textbf{Analysis.} This extends Problem 27. We reuse two ideas: (1) running in opposite directions, between two consecutive meetings the two together cover $1$ lap; (2) distance is proportional to speed.

\step{1}{Part a -- Divide the lap into 8 segments}
The ratio is $5:3$, so we divide the lap into $5+3=8$ equal segments and number the division points $0,1,2,\dots,7$ in Emma's running direction (point $0$ is $A$). Between two meetings the two together run $8$ segments, of which Emma runs $5$ and Noah $3$.\par
The $k$-th meeting is at the point $5k$ segments from $A$; every $8$ segments is a full lap, so we remove multiples of $8$:
\begin{center}\small
\begin{tabular}{|c|c|c|c|c|c|c|c|c|c|}
\hline
\rowcolor{tmBlue!10}Meeting $k$ & $1$ & $2$ & $3$ & $4$ & $5$ & $6$ & $7$ & $8$ & $9$\\ \hline
$5k$ & $5$ & $10$ & $15$ & $20$ & $25$ & $30$ & $35$ & $40$ & $45$\\ \hline
Point (remove multiples of $8$) & $5$ & $2$ & $7$ & $4$ & $1$ & $6$ & $3$ & $0$ & $5$\\ \hline
\end{tabular}
\end{center}
The first eight meetings fall on all $8$ different points $0,1,\dots,7$; the $9$th meeting is back at point $5$ and everything repeats. So there are \textbf{$8$ meeting points}.
\begin{figC}
\draw[Ink,line width=1.1pt] (0,0) circle (1.6);
\foreach \i/\k in {0/8,1/5,2/2,3/7,4/4,5/1,6/6,7/3}{
  \pgfmathsetmacro\rvabAng{90+45*\i}
  \fill[AnsRed] (\rvabAng:1.6) circle (2.6pt);
  \node[font=\scriptsize,anchor=\rvabAng+180,inner sep=2pt] at (\rvabAng:1.72) {point \i};
  \node[font=\scriptsize\bfseries,text=tmBlue] at (\rvabAng:1.25) {\k};}
\node[font=\scriptsize,text=tmBlue,align=center] at (0,0) {meeting\\no.};
\end{figC}
(Blue numbers inside: which meeting falls on that point; the points are numbered in Emma's running direction.)

\step{2}{Part b -- Same direction}
Running in the same direction, Emma overtakes Noah each time Emma has run exactly $1$ lap \emph{more} than Noah. In the same time Emma runs $5$ parts while Noah runs $3$, so Emma gains $2$ parts. Thus ``$1$ lap more'' corresponds to $2$ parts $=1$ lap, so each part is $\frac12$ lap, and Emma has run $5\cdot\frac12=2\frac12$ laps.\par
First overtaking: Emma has run $2\frac12$ laps, at the point opposite $A$. Second: $5$ laps, exactly at $A$. Then it repeats. So there are \textbf{$2$ points} (point $A$ and the point opposite $A$).

\step{3}{Part c -- The general rule}
Do the same as in part a with ratio $p:q$: divide the lap into $p+q$ segments; the $k$-th meeting is at the point $k\cdot p$ segments from $A$ (removing multiples of $p+q$).
\begin{proofbox}\textbf{If $p$ and $q$ are coprime, there are exactly $p+q$ meeting points.}\par
\emph{Meeting number $p+q$ is back at $A$:} because $(p+q)\cdot p$ is a multiple of $p+q$.\par
\emph{Meetings $1,2,\dots,p+q$ are at different points:} suppose meetings $i$ and $j$ ($i<j\le p+q$) are at the same point. Then $j\cdot p-i\cdot p=(j-i)\cdot p$ is a multiple of $p+q$. Any common divisor of $p$ and $p+q$ also divides $(p+q)-p=q$; since $p,q$ are coprime, $p$ and $p+q$ have no common prime factor. Hence every prime factor of $p+q$ must be in $j-i$, so $j-i$ is divisible by $p+q$. But $0<j-i<p+q$ -- impossible.\par
So there are exactly $p+q$ meeting points, equally spaced around the circle.\end{proofbox}
We need $p+q=10$ with $p>q\ge 1$ coprime:
\begin{center}
\begin{tabular}{|c|c|c|c|c|}
\hline
\rowcolor{tmBlue!10}$p:q$ & $9:1$ & $8:2$ & $7:3$ & $6:4$\\ \hline
Coprime? & yes & no ($=4:1$) & yes & no ($=3:2$)\\ \hline
Meeting points & $10$ & $5$ & $10$ & $5$\\ \hline
\end{tabular}
\end{center}
($5:5$ is excluded because $p>q$.) So there are two ratios: \textbf{$9:1$ and $7:3$}.
\begin{warnbox}The ratio $8:2$ has $8+2=10$ but is really $4:1$, which gives only $4+1=5$ meeting points. Always reduce the ratio before applying the ``$p+q$ points'' rule.\end{warnbox}
\verify{Check Problem 27 with the rule: the ratio $2:1$ gives $2+1=3$ points \checkmark. A computer simulation of all reduced ratios $p:q$ with $p\le 19$: only $9:1$ and $7:3$ give $10$ points.}
\ans{a) $8$ points;\ b) $2$ points;\ c) $9:1$ and $7:3$.}
}

% ------------------------------------------------------------------ 48
\begin{probbox}{48}{\stars{2}}
We call an \emph{increasing} list of positive integers a \emph{Fibo list} if, from the third number on, each number is the sum of the two numbers just before it (for example $11, 20, 31, 51, 82$ as in Problem 28).
\begin{enumerate}[label=\alph*),leftmargin=1.6em,itemsep=1pt]
\item Find all Fibo lists of $6$ numbers whose last number is $100$.
\item How many Fibo lists of $5$ numbers have last number $200$?
\item Prove that there is no Fibo list of $6$ numbers whose last number is $30$.
\end{enumerate}
\end{probbox}
\probsol{
\textbf{Analysis.} As in Problem 28, the whole list is determined by its first two numbers $a$ and $b$; the list is increasing exactly when $a<b$.

\step{1}{Express the numbers in terms of $a$ and $b$}
\begin{center}\small
\begin{tabular}{|c|c|c|c|c|c|}
\hline
\rowcolor{tmBlue!10}1st & 2nd & 3rd & 4th & 5th & 6th\\ \hline
$a$ & $b$ & $a+b$ & $a+2b$ & $2a+3b$ & $(a+2b)+(2a+3b)=3a+5b$\\ \hline
\end{tabular}
\end{center}
\begin{tipbox}The coefficients $1,1,2,3,5,8,\dots$ appearing in the table are the famous Fibonacci numbers!\end{tipbox}

\step{2}{Part a -- The equation $3a+5b=100$}
Conditions: $a\ge 1$, $a<b$, $3a+5b=100$.
\begin{itemize}[leftmargin=1.3em,itemsep=1pt]
\item From $a\ge 1$: $5b=100-3a\le 97$, so $b\le 19$.
\item From $a<b$: $100=3a+5b<3b+5b=8b$, so $b>12.5$, that is $b\ge 13$.
\item $3a=100-5b$ must be divisible by $3$.
\end{itemize}
Test $b$ from $13$ to $19$:
\begin{center}\small
\begin{tabular}{|c|c|c|c|c|c|c|c|}
\hline
\rowcolor{tmBlue!10}$b$ & $13$ & $14$ & $15$ & $16$ & $17$ & $18$ & $19$\\ \hline
$100-5b$ & $35$ & $30$ & $25$ & $20$ & $15$ & $10$ & $5$\\ \hline
$a=(100-5b)\div 3$ & $\times$ & $\mathbf{10}$ & $\times$ & $\times$ & $\mathbf{5}$ & $\times$ & $\times$\\ \hline
\end{tabular}
\end{center}
There are two lists: $\mathbf{10, 14, 24, 38, 62, 100}$ and $\mathbf{5, 17, 22, 39, 61, 100}$.

\step{3}{Part b -- The equation $2a+3b=200$}
Since $2a$ and $200$ are even, $3b$ is even, so $b$ is even: $b=2t$. Then $2a=200-6t$, that is $a=100-3t$.
\begin{itemize}[leftmargin=1.3em,itemsep=1pt]
\item $a\ge 1$: $100-3t\ge 1 \Rightarrow t\le 33$.
\item $a<b$: $100-3t<2t \Rightarrow 100<5t \Rightarrow t>20 \Rightarrow t\ge 21$.
\end{itemize}
So $t=21,22,\dots,33$: there are $33-21+1=13$ lists. (For example $t=21$: $37, 42, 79, 121, 200$; $t=33$: $1, 66, 67, 133, 200$.)
\begin{warnbox}For $t=20$ we get $a=b=40$: the list $40,40,80,120,200$ is \emph{not increasing} at the first step, so it is rejected. The number of integers from $21$ to $33$ is $33-21+1=13$, not $12$.\end{warnbox}

\step{4}{Part c -- No list ends in $30$}
We need $3a+5b=30$ with $1\le a<b$.
\begin{itemize}[leftmargin=1.3em,itemsep=1pt]
\item $3a=30-5b\ge 3$, so $5b\le 27$, that is $b\le 5$.
\item $5b=30-3a$ is divisible by $3$, so $b$ is divisible by $3$; with $b\le 5$ only $b=3$ is possible.
\item $b=3$ gives $3a=15$, $a=5$. But $a=5>3=b$ -- the list $5,3,8,11,19,30$ is not increasing.
\end{itemize}
So no Fibo list of $6$ numbers ends in $30$.
\begin{proofbox}A shorter argument: since $a<b$, the sixth number satisfies $3a+5b<3b+5b=8b$. To get $3a+5b=30$ we would need $b>30\div 8$, that is $b\ge 4$; together with $b\le 5$ and $b$ divisible by $3$, no value is left.\end{proofbox}
\verify{Testing every pair $1\le a<b$ with a computer: exactly $2$ lists in part a, $13$ lists in part b and $0$ lists in part c.}
\ans{a) $10,14,24,38,62,100$ and $5,17,22,39,61,100$;\ b) $13$ lists;\ c) see the solution.}
}

% ------------------------------------------------------------------ 49
\begin{probbox}{49}{\stars{1}}
Four friends Anna, Ben, Chloe and David each have at least one marble. We know that:
\begin{itemize}[leftmargin=1.6em,itemsep=0pt]
\item[(1)] Anna and Ben together have as many marbles as Chloe and David together.
\item[(2)] Anna and David together have more marbles than Ben and Chloe together.
\item[(3)] Ben has more marbles than Anna and Chloe together.
\end{itemize}
\begin{enumerate}[label=\alph*),leftmargin=1.6em,itemsep=1pt]
\item Put the four friends in order from fewest to most marbles.
\item Altogether the four friends have $30$ marbles. What is the greatest number of marbles Chloe can have? In that case, how many marbles does each friend have?
\end{enumerate}
\end{probbox}
\probsol{
\textbf{Analysis.} This extends Problem 29: instead of comparing pairs directly, the information here compares \emph{sums}. Let Anna, Ben, Chloe, David have $a, b, c, d$ marbles (positive integers). We have:
\[ \text{(1) } a+b=c+d,\qquad \text{(2) } a+d>b+c,\qquad \text{(3) } b>a+c. \]
\begin{recallbox}Adding (or subtracting) the two sides of an equation to (from) the corresponding sides of an inequality gives an inequality in the same direction: if $X>Y$ and $M=N$, then $X+M>Y+N$ and $X-M>Y-N$.\end{recallbox}

\step{1}{Compare Anna and Chloe: add (2) and (1)}
$(a+d)+(a+b)>(b+c)+(c+d)$, that is $2a+b+d>b+d+2c$. Subtract $b+d$ from both sides: $2a>2c$, so $a>c$: \textbf{Anna has more than Chloe}.

\step{2}{Compare Ben and David: subtract (1) from (2)}
$(a+d)-(a+b)>(b+c)-(c+d)$, that is $d-b>b-d$. Add $b+d$ to both sides: $2d>2b$, so $d>b$: \textbf{David has more than Ben}.

\step{3}{Compare Anna and Ben: use (3)}
$b>a+c\ge a+1>a$ (since $c\ge 1$). So \textbf{Ben has more than Anna}.

\step{4}{Part a -- Conclusion}
Putting the three results together: $c<a$, \ $a<b$, \ $b<d$, so $c<a<b<d$.\par
From fewest to most: \textbf{Chloe, Anna, Ben, David}.

\step{5}{Part b -- Use the total of 30}
From (1), $a+b=c+d$, and $(a+b)+(c+d)=30$, so $a+b=c+d=15$. Hence $b=15-a$.\par
Substitute into (3): $15-a>a+c$, that is $2a+c<15$; since $2a+c$ is a whole number, $2a+c\le 14$.\par
By part a, $a>c$, so $a\ge c+1$. Therefore $2(c+1)+c\le 2a+c\le 14$, that is $3c+2\le 14$, $3c\le 12$, $c\le 4$.

\step{6}{Part b -- Can $c=4$ actually happen?}
With $c=4$: we need $a\ge 5$ and $2a+4\le 14$, that is $a\le 5$. So $a=5$, $b=15-5=10$, $d=15-4=11$.\par
Check all three conditions:
(1) $5+10=15=4+11$ \checkmark;\quad (2) $5+11=16>14=10+4$ \checkmark;\quad (3) $10>5+4=9$ \checkmark.
\begin{warnbox}Proving ``$c\le 4$'' is only half the job: we must also show a concrete case with $c=4$ before we can claim that $4$ is the greatest value.\end{warnbox}
\verify{Listing with a computer every $(a,b,c,d)$ with total $30$ satisfying (1), (2), (3): there are $12$ of them, the largest $c$ is $4$, occurring only in $(5,10,4,11)$; and in every one of them $c<a<b<d$.}
\ans{a) Chloe, Anna, Ben, David;\ b) Chloe has at most $4$ marbles; then Anna has $5$, Ben $10$, Chloe $4$, David $11$.}
}

% ------------------------------------------------------------------ 50
\begin{probbox}{50}{\stars{1}}
A shop sells candy A at \$12 per kilogram and candy B at \$20 per kilogram.
\begin{enumerate}[label=\alph*),leftmargin=1.6em,itemsep=1pt]
\item How many grams of each kind must be mixed to get $1$~kg of mixture costing \$15 per kilogram (the mixture is worth exactly the total value of the candy put in)?
\item How many grams of candy C, costing \$30 per kilogram, must be added to the $1$~kg mixture of part a to get a new mixture costing \$18 per kilogram?
\item Prove that if candy A and candy B (both present) are mixed in any proportion, the price of the mixture is always more than \$12 and less than \$20 per kilogram.
\end{enumerate}
\end{probbox}
\probsol{
\textbf{Analysis.} This extends Problem 30. The principle: \emph{the total cost of the mixture $=$ the total cost of the parts mixed}. We measure mass in kilograms and money in dollars.
\begin{recallbox}Cost $=$ price (per $1$~kg) $\times$ mass (kg). Average price of a mixture $=$ total cost $\div$ total mass.\end{recallbox}

\step{1}{Part a -- Choose an unknown and write an equation}
Let the mass of candy A be $x$ (kg), with $0<x<1$; then the mass of candy B is $1-x$ (kg).\par
Cost of A: $12x$; cost of B: $20(1-x)=20-20x$; cost of $1$~kg of mixture: $15\cdot 1=15$ (dollars). So
\[ 12x+20-20x=15 \iff 20-8x=15 \iff 8x=5 \iff x=\frac{5}{8}=0.625. \]
So we need $0.625$~kg $=625$~g of candy A and $1000-625=375$~g of candy B.
\verify{$12\cdot 0.625+20\cdot 0.375=7.5+7.5=15$ (dollars) for $1$~kg, exactly \$15 per kilogram \checkmark.}

\step{2}{Part a -- Another way: compare with the price \$15 on a number line}
\begin{figC}
\draw[->,Navy,line width=1pt] (-0.5,0)--(5.4,0);
\foreach \x/\l in {0/12,1.8/15,4.8/20}{\draw[Navy,line width=1pt] (\x,0.12)--(\x,-0.12); \node[below,font=\small] at (\x,-0.12) {\l};}
\node[above,font=\small\bfseries,text=tmBlue] at (0,0.15) {A};
\node[above,font=\small\bfseries,text=AnsRed] at (1.8,0.15) {mixture};
\node[above,font=\small\bfseries,text=tmOrange!85!black] at (4.8,0.15) {B};
\node[anchor=west,font=\footnotesize,text=Slate] at (5.45,0) {dollars/kg};
\draw[Slate,dashed,line width=0.5pt] (0,-0.6)--(0,-1.05) (1.8,-0.6)--(1.8,-1.05) (4.8,-0.6)--(4.8,-1.05);
\draw[<->,tmBlue!72!black,line width=0.7pt] (0,-0.9)--(1.8,-0.9);
\draw[<->,tmBlue!72!black,line width=0.7pt] (1.8,-0.9)--(4.8,-0.9);
\node[below,font=\footnotesize] at (0.9,-0.95) {$15-12=3$};
\node[below,font=\footnotesize] at (3.3,-0.95) {$20-15=5$};
\end{figC}
Each kilogram of candy A is \$3 cheaper than \$15; each kilogram of candy B is \$5 dearer than \$15. The ``cheaper'' part must exactly balance the ``dearer'' part: $3x=5(1-x)$. Solving: $3x=5-5x$, $8x=5$, $x=0.625$ -- the same as Step 1. From $3x=5y$ (with $y=1-x$) we get $\dfrac{x}{y}=\dfrac{5}{3}$: the masses of A and B are in the ratio $5:3$, the two distances $3$ and $5$ \emph{reversed}.

\step{3}{Part b}
Let the mass of candy C added be $y$ (kg), $y>0$. Total cost: $15+30y$; total mass: $1+y$ (kg). For a price of \$18 per kilogram we need
\[ 15+30y=18(1+y)=18+18y \iff 30y-18y=18-15 \iff 12y=3 \iff y=0.25. \]
So we must add $0.25$~kg $=250$~g of candy C.
\verify{Total cost $15+30\cdot 0.25=15+7.5=22.5$ (dollars) for $1.25$~kg: $22.5\div 1.25=18$ \checkmark. (As in Step 2: $18$ is $3$ away from $15$ and $12$ away from $30$, so $3\cdot 1=12y$, $y=0.25$ \checkmark.)}

\step{4}{Part c}
Let the masses of A and B be $x$ and $y$ (kg), with $x>0$, $y>0$. The price of the mixture is $G=\dfrac{12x+20y}{x+y}$ (dollars per kg).
\begin{itemize}[leftmargin=1.3em,itemsep=1pt]
\item Since $y>0$, $20y>12y$, so $12x+20y>12x+12y=12(x+y)$. Dividing both sides by the positive number $x+y$: $G>12$.
\item Since $x>0$, $12x<20x$, so $12x+20y<20x+20y=20(x+y)$. Dividing both sides by $x+y$: $G<20$.
\end{itemize}
So $12<G<20$: the price of the mixture always lies \emph{between} the prices of the two candies.
\begin{warnbox}You cannot mix A and B to get a price of \$25 or of \$10 per kilogram: an average never goes outside the range from the lowest to the highest price.\end{warnbox}
\begin{parentbox}You can illustrate part c by mixing water: mixing water at $20^\circ$C with water at $60^\circ$C gives warm water in between, never hotter than $60^\circ$C. ``An average always lies in between'' is a very useful idea.\end{parentbox}
\ans{a) $625$~g of candy A and $375$~g of candy B;\ b) $250$~g of candy C;\ c) see the solution.}
}
